\begin{abstract}
LLM agents are rapidly being adopted for root-cause analysis (RCA) in microservice
systems, but their reported accuracy is difficult to trust: benchmark artifacts ---
fault-encoding run identifiers, injection containers named after their faults --- hand
the agent the answer. We present a leakage-controlled evaluation methodology (identifier
masking with cross-tool identity preservation, full-fidelity audit transcripts, and an
automated leak auditor) and apply it to a new dataset: 95 labeled fault injections
across two architecturally distinct microservice systems, recorded simultaneously in
\emph{four} time-aligned telemetry modalities --- metrics, logs, distributed traces,
and \emph{kernel traces} --- with pre-registered fault$\rightarrow$modality predictions.
We find that naming leaks alone inflate an agent's fault-classification accuracy by
57~percentage points; after closing them, a tool-using agent with generic
cross-modality tooling still reaches 83--87\% top-1 service localization, roughly
double two non-LLM baselines including published state of the art, at ${\sim}\$0.01$
per incident. Degradation experiments yield three counter-intuitive results: accuracy
survives $20\times$ trace sampling unchanged; removing kernel telemetry entirely costs
little accuracy but 34\% more investigation effort; and a \emph{partial} kernel view
(raw KPIs without interpretive framing) is \emph{worse than none}. An experience layer
of reusable ``skills'' helps only when skill selection is reliable --- mis-selected
skills cost more than they contribute. All 500+ diagnoses ship with machine-checkable
no-leakage certificates. \todo{tighten to venue word limit}
\end{abstract}

%% TODO: CCS concepts + keywords (required by acmart)
\begin{CCSXML}
<ccs2012>
<concept><concept_id>10011007.10011074.10011099</concept_id>
<concept_desc>Software and its engineering~Software verification and validation</concept_desc>
<concept_significance>500</concept_significance></concept>
</ccs2012>
\end{CCSXML}
\ccsdesc[500]{Software and its engineering~Software verification and validation}
\keywords{root-cause analysis, LLM agents, observability, kernel tracing, benchmark
leakage, microservices, AIOps}
